% Generated draft: requires mathematical and source review.
\documentclass[11pt]{article}
\usepackage[margin=25mm]{geometry}
\usepackage{fontspec}
\setmainfont{Latin Modern Roman}
\usepackage{amsmath,amssymb,mathtools,unicode-math}
\setmathfont{Latin Modern Math}
\usepackage{enumitem,needspace,xurl,hyperref}
\setlength{\parindent}{0pt}
\setlength{\parskip}{5pt}
\setlength{\emergencystretch}{3em}
\newcommand{\N}{\mathbb{N}}
\newcommand{\Z}{\mathbb{Z}}
\newcommand{\Q}{\mathbb{Q}}
\newcommand{\R}{\mathbb{R}}
\newcommand{\C}{\mathbb{C}}
\newcommand{\F}{\mathbb{F}}
\newcommand{\A}{\mathbb{A}}
\newcommand{\PP}{\mathbb{P}}
\newcommand{\E}{\mathbb{E}}
\newcommand{\eps}{\varepsilon}
\newcommand{\dd}{\,\mathrm{d}}
\begin{document}
\section*{UnsolvedMath ID: 30001784}
\textbf{Source problem code:} catalogue entry
\subsection{Definitions and mathematical statement}
For terminology, a translation surface is a compact connected surface equipped, away from finitely many conical singularities, with coordinate charts whose transition maps are translations. Its Veech group is the subgroup of \(\mathrm{SL}_2(\R)\) consisting of derivatives of orientation-preserving affine automorphisms. A Veech surface is a translation surface whose Veech group is a lattice. In this setting, non-arithmetic means that the Veech group is not conjugate to a group commensurable with \(\mathrm{SL}_2(\Z)\); two groups are commensurable if their intersection has finite index in each.

For a specified translation surface \(Y\), let \(\phi_t^\theta\) denote its straight-line flow in direction \(\theta\), defined almost everywhere with respect to flat area. Write
\[
N(Y)=\{\theta\in\R/\pi\Z:\phi_t^\theta\text{ is not ergodic with respect to normalized flat area}\}.
\]
Ergodicity means that every measurable invariant set has measure zero or one. For a specified ambient moduli space \(\mathcal M\) and associated Teichm\"uller geodesic \(\gamma_{Y,\theta}:[0,\infty)\to\mathcal M\), one may define
\[
D_{\mathcal M}(Y)=\{\theta\in\R/\pi\Z:\gamma_{Y,\theta}\text{ eventually leaves every compact subset of }\mathcal M\}.
\]
Hausdorff dimensions of direction sets are taken using the usual angular metric on \(\R/\pi\Z\). These definitions explain conventional terminology; the supplied record does not identify whether these are precisely the objects intended by its question.

\textbf{Source question, with unresolved qualifications.} Does the Hausdorff-dimension dichotomy for non-ergodic or divergent directions persist when the square torus \(\R^2/\Z^2\) is replaced by a non-arithmetic Veech surface, and what is the analogue of Perez Marco's condition?

The supplied record does not state the original dimension dichotomy, its numerical alternatives, its hypotheses, or Perez Marco's condition. It also does not identify whether the relevant directional flows are on the replacement surface itself or on a construction over that surface, nor does it specify the ambient space defining divergence. Consequently, the displayed sets are explanatory notation only: the record does not determine a fully quantified assertion about their dimensions. The intended family, parameter domains, exceptional cases, and quantifiers must be recovered from the cited report before a precise conjecture can be stated.
\subsection{Short English statement}
The question asks whether a Hausdorff-dimension dichotomy associated with the square torus extends to non-arithmetic Veech surfaces, and whether Perez Marco's condition has an analogue there. The supplied record omits the dichotomy and the condition themselves, together with essential details of the underlying construction and divergence space. This is therefore a qualified statement-only draft requiring human mathematical review.
\subsection{Sources}
\begin{itemize}
\item Supplied local UnsolvedMath registry record, permanent ID \(30001784\), title \emph{Direction-Dimension Dichotomies on Nonarithmetic Veech Surfaces}, exact source code \texttt{catalogue entry}, in the \emph{Oberwolfach Reports Open Problems} collection. The record gives proposed year 2011 and identifies no individual proposer.
\item \emph{Billiards, Flat Surfaces, and Dynamics on Moduli Spaces} (2011), Oberwolfach Reports, pp.\ 1361--1427, DOI \texttt{10.4171/owr/2011/25}: \url{https://doi.org/10.4171/owr/2011/25}. Supplied literature-source label: ``original OWR report.'' The record describes Oberwolfach Reports as documenting mathematical workshops and their open-problem discussions.
\item Supplied literature-source label: ``weak-mixing results,'' \url{https://doi.org/10.1090/jams/856}. No further bibliographic metadata is supplied for this reference.
\item Literature assessment dated \texttt{2026-08-21}: the supplied curation labels the problem open and states that the extension from arithmetic to non-arithmetic Veech surfaces remains a question, while later weak-mixing results do not establish the requested non-ergodic or divergent dimension dichotomy or a general Perez--Marco analogue. This is supplied dated curation, not an independently verified assessment; its time-sensitive status and cited claims require independent checking.
\item The registry labels the statement ``exact'' and reports that its mathematical content was checked against the official Oberwolfach report, with notation, context, and question form normalized without changing the source-level claim. This verification claim is supplied metadata and has not been independently verified here. The record names Perez Marco's condition but supplies neither its statement nor a separate reference for it.
\end{itemize}
\end{document}
